\section*{Bab \ref{ch:b3:field-lab-methods-3} --- Metode Lapangan dan Laboratorium III}
% Values from the Book 5 scratch calculator calc/ch31.py.

\begin{solution}{exo:b3:field-lab-methods-3:1}
(a)~$\qty{0.01}{\milli\gal} = \qty{e-7}{m.s^{-2}}$, dan
$\num{e-7}/\num{9.81} = \num{1.0e-8}$: satu bagian dalam seratus juta.

(b)~$0.01/0.3086 = \qty{0.032}{m}$, sekitar \qty{3.2}{cm}.

(c)~Bersama-sama kedua efek itu mengubah gravitasi sebesar
\qty{0.1967}{\milli\gal.m^{-1}}, selisihnya $0.3086 - 0.1119$, sehingga
tinggi harus berubah sebesar $0.01/0.1967 = \qty{5.1}{cm}$.
\end{solution}

\begin{solution}{exo:b3:field-lab-methods-3:2}
$\Delta m = \tfrac43\pi(\qty{50}{m})^3\times\qty{500}{kg.m^{-3}} =
\qty{2.62e8}{kg}$. Puncaknya adalah $G\Delta m/z^2 =
\num{6.674e-11}\times\num{2.62e8}/\num{e4} = \qty{1.75e-6}{m.s^{-2}}$,
yaitu \qty{0.175}{\milli\gal}. Setengah lebarnya adalah $x_{1/2} = z/1.305 =
\qty{77}{m}$.
\end{solution}

\begin{solution}{exo:b3:field-lab-methods-3:3}
(a)~Dengan $AM = MN = NB = a$, $K = 2\pi/(1/a - 1/2a - 1/2a + 1/a) =
2\pi a = \qty{62.83}{m}$, dan $\rho_a = 62.83\times0.239/0.1 =
\qty{150}{\ohm\metre}$. (b)~$K = \pi(15^2 - 5^2)/(2\times5) = \pi\times
200/10 = \qty{62.83}{m}$. Elektrodanya kemudian berada pada $0$, $10$, $20$
dan \qty{30}{m}, dalam urutan AMNB: ini adalah susunan Wenner dengan spasi
\qty{10}{m}.
\end{solution}

\begin{solution}{exo:b3:field-lab-methods-3:4}
(a)~$V = 0.2998/\sqrt5 = \qty{0.134}{m.ns^{-1}}$. (b)~$t =
2\times10/0.134 = \qty{149}{ns}$. (c)~$\delta = 0.0053\sqrt5/\num{e-3} =
\qty{11.9}{m}$, dan kedalaman berguna sekitar \qty{36}{m}. (d)~Di lempung
$\delta = 0.0053\sqrt{20}/0.01 = \qty{2.4}{m}$, kedalaman berguna sekitar
\qty{7}{m}, dan kecepatannya hanya \qty{0.067}{m.ns^{-1}}. Pendekatan
masih berlaku di sana: pada \qty{100}{MHz},
$\omega\varepsilon = 2\pi\times\num{e8}\times20\times\num{8.85e-12} =
\qty{0.11}{S.m^{-1}}$, sepuluh kali $\sigma$.
\end{solution}

\begin{solution}{exo:b3:field-lab-methods-3:5}
(a)~Silinder memiliki $\lambda = \pi R^2\Delta\rho$ per satuan panjang, dan
di luarnya tarikannya sama dengan tarikan massa garis pada sumbunya (teorema Gauss
pada silinder koaksial): $\Delta g = 2G\pi R^2\Delta\rho\,z/(x^2
+ z^2)$. Setengah puncak tercapai saat $x^2 + z^2 = 2z^2$, sehingga $z =
x_{1/2}$. (b)~Bola berada pada $1.305\times100 = \qty{130}{m}$, sedangkan
silinder pada \qty{100}{m}: profil yang sama, jika dibaca sebagai bola, menempatkan
benda \qty{30}{\percent} lebih dalam. Anomali bola meluruh lebih cepat seiring
jarak, sehingga bola harus lebih dalam agar sama lebarnya.
\end{solution}

\begin{solution}{exo:b3:field-lab-methods-3:6}
(a)~Setiap titik mewakili $\qty{100}{m}\times\qty{100}{m} =
\qty{e4}{m^2}$, sehingga $\iint\Delta g\,\mathrm{d}A \approx
\num{-438e-5}\times\num{e4} = \qty{-43.8}{m^3.s^{-2}}$, dan $\Delta m
= -43.8/(2\pi\times\num{6.674e-11}) = \qty{-1.04e11}{kg}$. (b)~$1.04/1.74
= \qty{60}{\percent}$. Sisanya ada di lereng anomali di luar
peta: granit mencapai kedalaman lebih dari satu kilometer, dan
anomali bagian dalamnya meluruh perlahan terhadap jarak. Memperluas peta
cukup jauh, atau memakai model, akan memulihkannya.
\end{solution}

\begin{solution}{exo:b3:field-lab-methods-3:7}
(a)~Untuk retas tipis, anomali turun menjadi setengah puncaknya pada $x = \pm z$:
lebar penuhnya pada setengah puncak adalah $2z = \qty{60}{m}$, dan tiga
pembacaan melintasinya memerlukan spasi paling besar \qty{20}{m}. (b)~$2\times
220/3 = \qty{147}{m}$. (c)~Spasi stasiun yang diperlukan bergantung pada
kedalaman benda terdangkal yang diminati: survei yang dirancang untuk retas
dalam, pada \qty{150}{m}, akan menimbulkan aliasing pada anomali benda dangkal, yang
hanya \qty{60}{m} lebarnya; survei yang dirancang untuk benda dangkal juga melihat
benda dalam itu juga, dengan pembacaan jauh lebih banyak daripada yang diperlukan.
\end{solution}

\begin{solution}{exo:b3:field-lab-methods-3:8}
(a)~$F = 80/8 = 10$ dan $\phi = F^{-1/m} = 10^{-1/2} = \num{0.316}$.
(b)~$10^{-1/1.8} = \num{0.278}$ dan $10^{-1/2.2} = \num{0.351}$. Perubahan
sebesar \num{0.2} pada $m$ menggeser porositas sekitar empat poin persentase;
resistivitas, yang terukur hingga beberapa persen, menggesernya jauh lebih kecil. Eksponen
sementasi mengendalikan ketidakpastian, dan ia harus dikalibrasi
pada sampel batuan yang sama.
\end{solution}

\begin{solution}{exo:b3:field-lab-methods-3:9}
(a)~Pada $s = 2h_1$, $(2nh_1/s)^2 = n^2$, sehingga sukunya adalah $k^n(1 +
n^2)^{-3/2}$: $0.3536k$, $0.0894k^2$, $0.0316k^3$. Untuk $k = 0.5$:
$\num{0.1768}$, $\num{0.0224}$, $\num{0.0040}$, dan $\rho_a/\rho_1 \approx
1 + 2\times0.2031 = \num{1.406}$ (jumlah penuh memberikan \num{1.409}). Untuk $k
= -0.5$ suku-sukunya bergantian tanda, $-0.1768 + 0.0224 - 0.0040$, dan
$\rho_a/\rho_1 \approx \num{0.683}$ (jumlah penuh \num{0.685}). (b)~Untuk $s \ll
h_1$ setiap suku menghilang dan $\rho_a \to \rho_1$; untuk $s \gg h_1$ setiap
kurung mendekati~1 dan $1 + 2\sum k^n = 1 + 2k/(1 - k) = (1 + k)/(1 -
k) = \rho_2/\rho_1$. (c)~$\rho_2/\rho_1 = (1 + k)/(1 - k) = 3$ dan $1/3$.
\end{solution}

\begin{solution}{exo:b3:field-lab-methods-3:10}
(a)~Dalam urutan AMNB dengan spasi $a$: $AM = a$, $BM = 2a$, $AN = 2a$,
$BN = a$, sehingga kurungnya adalah $2/a - 1/a = 1/a$ dan $K = 2\pi a$.
(b)~Dengan A dan B pada $\mp s$ dan M dan N pada $\mp l$: $AM = BN = s - l$,
$BM = AN = s + l$, kurungnya adalah $2/(s - l) - 2/(s + l) = 4l/(s^2 -
l^2)$, dan $K = \pi(s^2 - l^2)/(2l)$. Ketika $l \ll s$, $\Delta V_{MN}
\approx 2l\,E$ dengan $E = 2(-\mathrm{d}V/\mathrm{d}r)$ di pusat, medan
kedua elektroda arus; lalu $\rho_a = K\Delta V/I \to
\pi s^2/(2l)\times 4l(-\mathrm{d}V/\mathrm{d}r)/I = 2\pi s^2(-
\mathrm{d}V/\mathrm{d}r)/I$. (c)~Rasio faktor pendekatan terhadap
faktor eksaknya adalah $s^2/(s^2 - l^2) = 1/(1 - 0.01) = 1.010$: satu persen.
\end{solution}

\begin{solution}{exo:b3:field-lab-methods-3:11}
(a)~Pembacaan: \qty{0.010}{\milli\gal}. Tinggi:
$0.1967\times0.05 = \qty{0.010}{\milli\gal}$. Massa jenis: koreksi Bouguer
adalah $0.04191\,h\,d$ dengan $d$ dalam \unit{g.cm^{-3}}, sehingga
$0.04191\times50\times0.1 = \qty{0.210}{\milli\gal}$. \omterm{def:b3:field-lab-methods-3:terrain}{Koreksi medan}:
\qty{0.200}{\milli\gal}. Total: $\sqrt{0.010^2 + 0.010^2 + 0.210^2 +
0.200^2} = \qty{0.29}{\milli\gal}$. (b)~Massa jenis dan \omterm{def:b3:field-lab-methods-3:terrain}{koreksi medan}
mendominasi, masing-masing sekitar dua puluh kali kesalahan pembacaan. Untuk mendeteksi
\qty{0.1}{\milli\gal}, massa jenis lempengan harus diukur pada
sampel atau dicocokkan dengan data (nilai yang paling sedikit mencerminkan
topografi), dan medannya dikoreksi dari model elevasi digital yang rinci.
(c)~$\sqrt{0.010^2 + 0.010^2 + 0.210^2} = \qty{0.21}{\milli\gal}$.
\end{solution}

\begin{solution}{exo:b3:field-lab-methods-3:12}
(a)~Massa yang sesuai dengan pembacaan pusat adalah $\Delta m = \Delta g(0)z^2/G
= \num{-4.16e-5}z^2/\num{6.674e-11}$: \qty{-5.61e10}{kg} pada
\qty{300}{m}, \qty{-9.97e10}{kg} pada \qty{400}{m}, \qty{-1.56e11}{kg} pada
\qty{500}{m}. Pembacaan prediksi pada $x = \pm300$ dan $\pm600$ adalah
\num{-1.47} dan \num{-0.37}, \num{-2.13} dan \num{-0.71}, \num{-2.62}
dan \qty{-1.09}{\milli\gal}. \omterm{def:b3:field-lab-methods-3:inversion-practice}{Ketakcocokannya} adalah $\chi^2 = 1238$, $3.7$ dan
$856$. (b)~Hanya model pada \qty{400}{m}, dengan $\chi^2 = 3.7$ untuk lima
data, berada dalam batas galat. (c)~Ya: setiap benda simetris bola dengan massa dan pusat yang sama
memberikan anomali yang persis sama, dan bentuk lain
dapat berada dalam batas galat. Massa jenis yang diukur pada sampel, sebuah
sumur bor atau metode lain akan memutuskan.
\end{solution}

\begin{solution}{pb:b3:field-lab-methods-3:1}
\textbf{1.} $(2846.708 - 2846.512)/\qty{2.0}{h} = \qty{0.098}{\milli\gal.h^{-1}}$.

\textbf{2.} $2851.075 - 0.098\times0.5 - 2846.512 = \qty{4.514}{\milli\gal}$;
$2850.497 - 0.088 - 2846.512 = \qty{3.897}{\milli\gal}$;
$2849.392 - 0.127 - 2846.512 = \qty{2.753}{\milli\gal}$.

\textbf{3.} Ketinggian di atas dasar adalah $-27.1$, $-30.6$ dan
\qty{-18.1}{m}; dengan menambahkan $0.1967\,\Delta h$ diperoleh anomali Bouguer
$\num{-0.817}$, $\num{-2.122}$ dan \qty{-0.808}{\milli\gal}. Stasiun
berada lebih rendah daripada dasar dan lebih dekat ke pusat Bumi,
sehingga pembacaan mentahnya lebih tinggi; koreksi menghilangkan hal itu dan
menyisakan anomali rendah di atas granit.

\textbf{4.} Selisih tambahan $0.098 - 0.040 = \qty{0.058}{\milli\gal.h^{-1}}$,
\qty{0.12}{\milli\gal} selama satu lintasan tertutup, adalah pasang surut Bumi padat. Selama dua
jam ia berubah hampir linear, sehingga garis lurus yang melewati dua
pembacaan dasar menghilangkannya bersama hanyutan.

\textbf{5.} Galatnya adalah $\num{0.008}$, $\num{0.009}$ dan
\qty{0.017}{\milli\gal}. Pembulatan setiap ketinggian ke \qty{0.1}{m} dapat membuat
selisih ketinggian keliru hingga \qty{0.1}{m}, yaitu
\qty{0.020}{\milli\gal}: pembulatan menjelaskan galat tersebut, sedangkan pasang surut
yang tersisa setelah garis lurus berada di bawah \qty{0.001}{\milli\gal}.

\textbf{6.} $0.01/0.1967 = \qty{5.1}{cm}$. \qty{3.3}{cm} pada manual adalah
$0.01/0.3086$: anomali udara-bebas tidak bergantung pada massa jenis
yang dipilih untuk koreksi Bouguer, dan ketinggian juga harus mengikutinya.

\textbf{7.} $x_{1/2} = (3.952 - 3.048)/2 = \qty{0.452}{km}$, dan $z = 1.305
\times452 = \qty{590}{m}$ di bawah permukaan tanah.

\textbf{8.} $\Delta m = \Delta g\,z^2/G = \num{-2.293e-5}\times590^2/
\num{6.674e-11} = \qty{-1.20e11}{kg}$. Dengan $\Delta\rho =
\qty{-130}{kg.m^{-3}}$ volumenya adalah \qty{9.2e8}{m^3} dan jari-jarinya
\qty{603}{m}, lebih besar daripada kedalaman pusatnya: bola itu akan
menembus permukaan tanah. Tidak ada bola dengan kontras massa jenis granit
yang sesuai; benda itu harus mencapai dekat permukaan.

\textbf{9.} Granit membentang dari permukaan hingga kedalaman besar. Bagian dangkalnya
di dekat stasiun membuat puncaknya sempit, dan aturan itu, yang
hanya mengetahui lebarnya, menempatkan satu bola pada kedalaman sedang; bagian
dalam menambahkan sisi lebar yang tidak dimiliki bola (lihat gambar
profil), dan massa yang mereka kandung terlewatkan.

\textbf{10.} $2\pi G\Delta\rho = 2\pi\times\num{6.674e-11}\times(-130) =
\qty{-5.45e-8}{s^{-2}}$. Dengan $z_1 = 0$, $z_2 = \qty{1645}{m}$, $a =
\qty{461}{m}$: nilai dalam kurung adalah $1645 + 461 - 1708.4 = \qty{397.6}{m}$,
dan $\Delta g = \qty{-2.17}{\milli\gal}$. Untuk kedalaman tak terbatas, nilai
dalam kurung cenderung ke $a$, dan $\Delta g = \qty{-2.51}{\milli\gal}$.

\textbf{11.} Nilai dalam kurung harus $2.29/\num{5.45e-3} = \qty{420}{m}$, yang
dapat dicapai dengan dasar sekitar \qty{2.6}{km}. Dasar apa pun antara sekitar
\qty{1.6}{km} dan kedalaman tak terbatas hanya mengubah puncak beberapa persepuluh
miligal saja, dan benda yang sedikit lebih lebar atau lebih padat memberikan puncak yang sama
dengan dasar yang lebih dangkal: dasar itu tidak dapat diresolusi.

\textbf{12.} Dengan $z_1 = \qty{6}{m}$: $1639 + 461.0 - 1708.4 = 391.6$, dan
$\Delta g = \qty{-2.14}{\milli\gal}$, \qty{0.03}{\milli\gal} lebih kecil daripada
dengan puncak di permukaan tanah. Selisih ini berada di bawah galat medan dan massa jenis
survei: gravitasi tidak dapat membedakan apakah granit mencapai
permukaan di bawah aluvium. Sounding listrik dapat.

\textbf{13.} $K = \pi(20^2 - 1^2)/(2\times1) = \qty{626.7}{m}$, dan $\Delta V =
\rho_a I/K = 821\times0.050/626.7 = \qty{65.5}{mV}$.

\textbf{14.} $V = 2\times1/\sqrt{17.1^2 - 10.7^2} = 2/13.34 =
\qty{0.150}{m.ns^{-1}}$, $\varepsilon_r = (0.2998/0.150)^2 = 4.0$, yaitu
pasir kering; pipa tersebut berada pada kedalaman $0.150\times10.7/2 = \qty{0.80}{m}$.

\textbf{15.} $0.150\times20.0/2 = \qty{1.50}{m}$.

\textbf{16.} $20/821 = \num{0.0244}$, $30/1309 = \num{0.0229}$, $40/1684 =
\num{0.0238}$: $S \approx \qty{0.0237}{S}$.

\textbf{17.} Lapisan kering menyumbang $1.5/600 = \qty{0.0025}{S}$, sedangkan
lapisan jenuh $\qty{0.0212}{S}$, sehingga tebalnya $0.0212\times222 =
\qty{4.7}{m}$, dan granit berada $1.5 + 4.7 = \qty{6.2}{m}$ di bawah
dasar. Taksiran ini tinggi: granit tidak sepenuhnya resistif tak terhingga, sehingga
pada spasi ini $\rho_a$ masih berada di bawah garis
$s/S$, dan $s/\rho_a$ melebihkan $S$. (Model daerah itu memiliki
\qty{6}{m} aluvium.)

\textbf{18.} $\phi = (20/222)^{1/2} = \num{0.30}$. Dengan $m = 1.3$, $\phi =
0.0901^{1/1.3} = \num{0.16}$; dengan $m = 2.8$, $\phi = \num{0.42}$.

\textbf{19.} Dengan mengurangkan nilai dasar dari lintasan, diperoleh dari barat ke timur:
3.4, 4.1, 4.8, 5.9, 7.4, 9.2, 11.5, 14.2, 17.2, 19.8, 21.0, 20.1, 17.6,
14.5, 11.5, 9.1, 7.2, 5.7, 4.6, 3.8 dan \qty{3.1}{nT}.

\textbf{20.} Tanpa koreksi, puncaknya \qty{28.0}{nT} pada \qty{2.20}{km}, dan
setengahnya, \qty{14.0}{nT}, tidak pernah dicapai ke arah timur: lintasan berakhir
pada \qty{15.2}{nT}, karena medan naik \qty{12}{nT} selama
pengukuran.

\textbf{21.} Puncaknya \qty{21.0}{nT} pada \qty{2.20}{km}; setengahnya adalah
\qty{10.5}{nT}. Barat: antara \num{9.2} pada \num{1.95} dan \num{11.5} pada
\qty{2.00}{km}, pada $1.95 + 0.05\times1.3/2.3 = \qty{1.978}{km}$. Timur:
antara \num{11.5} pada \num{2.40} dan \num{9.1} pada \qty{2.45}{km}, pada
$2.40 + 0.05\times1.0/2.4 = \qty{2.421}{km}$. Lebar setengahnya adalah $(2.421
- 1.978)/2 = \qty{0.221}{km}$.

\textbf{22.} Anomali simetris dalam medan yang miring \ang{65} berarti hanya
magnetisasi vertikal yang bekerja: benda itu berjurus sepanjang meridian magnetik,
utara--selatan. Aturan dari
\cref{prop:b3:field-lab-methods-3:dyke} mengasumsikan retas yang tipis, vertikal,
dua dimensi, yang menjangkau jauh ke bawah, termagnetisasi secara induksi, dengan puncak
yang berperan sebagai garis muatan magnetik.

\textbf{23.} $\kappa t = 2\pi z\,\Delta T_\mathrm{max}/(F\sin^2 I) = 2\pi\times
221\times21.0/(48000\times0.8214) = \qty{0.74}{m}$. Dengan $\kappa =
0.03$, $t = \qty{25}{m}$. Hanya hasil kalinya yang tetap: retas dua kali
lebih tebal dengan suseptibilitas setengahnya memberikan anomali yang sama.

\textbf{24.} Anomali itu selebar \qty{443}{m} pada setengah puncaknya, dan stasiunnya
berjarak \qty{50}{m}: sembilan pembacaan melintasinya, tiga kali lipat
minimum aturan itu.

\textbf{25.} Puncak retas berada
\textbf{\qty[text-series-to-math=true]{220}{m}} di bawah permukaan tanah. Model
geologi meletakkan ketidakselarasan sedalam \qty{222}{m} pada titik ini:
retas naik hingga dasar penutup dan terpotong di sana, lebih tua
daripada batu pasir Trias yang menimbunnya.
\end{solution}
